\section{補助空間の22既約成分：帰属の完結}\label{sec:Q22}
ここまでの結果と、v3--v6 の固定モジュラー計算を合わせると、
補助空間の被約既約成分の分類を完結できる。
トレース $s=Q^{12}+Q^{34}$ に対し、互換条件の下で
\[
 \delta Q=X_s=(-s_y,s_x,-s_w,s_z),\qquad
 \delta\Jac(Q)=\mathcal L_{\delta Q}Q
\]
である。従って固定 $s$ における $X_s$ 不変性の線形方程式は Jacobi ideal に含まれる。
不変部分だけへ移る際に、解や微小構造を勝手に捨てたわけではない。

\begin{theorem}\label{thm:Q22}
$\cQ_{\mathrm{red}}$ は正確に22個の既約成分を持つ。
その分類と次元は次表の通りである。
\end{theorem}
\begin{center}\small
\begin{tabular}{p{0.26\linewidth}p{0.30\linewidth}cc}
\toprule
一般モジュラー型&成分の名称&個数&各成分の affine dimension\\\midrule
正則半単純、非共鳴&$G_1,G_2$&2&11\\
正則半単純、共鳴 $1:3$&$T_1,T_2,T_3,T_4$&4&10\\
非正則半単純&$E_1,E_2,E_3,E_4$&4&10\\
正則冪零 $(4)$&$N_B$&1&9\\
最小冪零 $(2,1,1)$&$M_0,M_1,M$&3&$13,13,10$\\
冪零 $(2,2)$&開枝、$C_+,C_-,D_\pm,U_1,A_0,AB$&7&10\\
零&$\mathcal L_0$&1&10\\\bottomrule
\end{tabular}
\end{center}
特に次元別には、13次元が2個、11次元が2個、10次元が17個、9次元が1個である。
\begin{proof}
非冪零モジュラー場を持つ全点は、v5 の完全な開集合の分類により最初の10成分に含まれる。
非半単純かつ非冪零の二つの混合 Jordan 型もその際に処理済みである。
一方の固定 fiber は空で、他方の三つの直線は実際の多項式変形により $G_1,G_2$ の境界へ入る。

$\mathfrak{sp}_4$ の非零冪零軌道は $(4),(2,2),(2,1,1)$ の三つである。
正則冪零型では固定 fiber の二直線の一方が既存成分の境界、他方が9次元の全成分 $N_B$ を与える。
最小冪零型では9、9、6次元の固定 fiber 成分の飽和が、13、13、10次元の三成分となり、
最後の帰属は定理\ref{thm:minimal-full}で閉じた。
$(2,2)$ 型の固定 fiber の17成分は、中心化群による同一視後に11飽和となる。
本ノートでそのうち四つが $G_1,G_2$ の境界、残り七つが相異なる全成分であると決定した。

最後に零 fiber は定理\ref{thm:zero-three-components}の三成分で尽くされる。
$\mathcal R_0$ は既出の13次元成分、$\mathcal C_0$ は $G_1\cap G_2$ に含まれ、
$\mathcal L_0$ 自身が10次元の全成分である。
従って全点が表の22個の閉成分に入る。
各成分の全成分性と相異なることは、既出の接空間・形式局所環・固定 fiber の中心化群作用で証明済みである。
有限個のこれらの閉集合が全点を覆うので、さらに別の既約成分は存在しない。
\end{proof}

\subsection{代表式を参照するための索引}
$G_1,G_2$ は $(2,2)$ 型の節に記した $N_1,N_2$ の飽和閉包である。
正則半単純の四成分には、$u=xy,v=zw$ として $N=N(A,B,C,D,E;P,R)$ を
\begin{align*}
 N^{12}&=Au+Bv,& N^{34}&=Cu+Dv,\\
 N^{13}&=Exz,&N^{14}&=(C-E)xw+Py^2,\\
 N^{23}&=(-B-E)yz+Rx^2,&N^{24}&=(B-C+E)yw
\end{align*}
で定め、次の値における $\Sp_4\times\CC^*$ 軌道閉包を取る。
\begin{center}
\begin{tabular}{c|rrrrrrr}
 & $A$&$B$&$C$&$D$&$E$&$P$&$R$\\\hline
$T_1$&1&0&0&3&1&1&0\\
$T_2$&$1/2$&0&$1/2$&3&1&1&0\\
$T_3$&1&$3/2$&0&$3/2$&$-1/2$&1&0\\
$T_4$&$-1$&0&2&3&1&1&3
\end{tabular}
\end{center}
非正則半単純の四代表は、成分順序 $(12,13,14,23,24,34)$ で
\begin{align*}
 E_1&=(xy,xz,-2xy-xw,-yz,y^2+yw,zw),\\
 E_2&=(xy/2,xz,-2xy-xw/2,-yz,y^2+yw/2,xy/2+zw),\\
 E_3&=(xy,xz/2,-xw/2+yz,-yz/2,yw/2,zw),\\
 E_4&=(xy/2,xz/2,yz,-yz/2,0,xy/2+zw)
\end{align*}
である。正則冪零の $N_B$ は
\[
 (3xw/2,-3x^2/2,3xz/2,3xy/2-zw/2,-3yz/2+w^2/2,-xw/2+z^2/2)
\]
の $\Sp_4$ 軌道閉包を取る。

最小冪零の既存二成分 $M_0,M_1$ は、固定トレース $x^2/2$ の次の9パラメータ族を飽和して得る。
$c=0$ が $M_0$、$c=1/4$ が $M_1$ に対応する。
\begin{align*}
 Q^{12}&=(1/2-c)x^2,\qquad Q^{13}=Q^{14}=0,\qquad Q^{34}=cx^2,\\
 Q^{23}&=t_1x^2+(-2c-t_6)xz-\tfrac12t_8z^2-2t_9zw+t_2xw+t_3w^2,\\
 Q^{24}&=t_4x^2+t_5xz+t_6xw+t_7z^2+t_8zw+t_9w^2.
\end{align*}
新成分 $M$ の代表は前節の $Q_M$ である。
$(2,2)$ 型の七成分の代表と全パラメータは \path{data/nil22_classification_summary.json} と
当該節の共通正規形に記した。$\mathcal L_0$ は零 fiber の節の二パラメータ族の飽和閉包である。
これらは成分を決定する代表族であり、全点の群軌道に対する一意な標準形を主張しているわけではない。

\subsection{一般横断環が示す三つの異なる厚み}
今回の一般非被約構造には
\[
 K[\varepsilon]/(\varepsilon^3),\qquad
 K[\varepsilon]/(\varepsilon^4),\qquad
 K[u,v]/(u,v)^2
\]
が現れる。最初と最後は長さが同じ3でも同型でない。
最大イデアルを $\mathfrak m$ とすると、前者は
$\dim_K\mathfrak m/\mathfrak m^2=1$、$\mathfrak m^2\ne0$、
後者は $\dim_K\mathfrak m/\mathfrak m^2=2$、$\mathfrak m^2=0$ である。
本ノートの22成分分類は被約台の分類であり、全境界における primary decomposition、
embedded components、成分交叉の局所環をすべて求めたわけではない。
